\documentclass[10pt,a4paper]{article}
\usepackage[utf8]{inputenc}
\usepackage[T1]{fontenc}
\usepackage[spanish,es-noshorthands]{babel}
\usepackage[margin=3cm]{geometry}
\usepackage{amsmath,amssymb}
\usepackage{booktabs,tabularx,array}
\usepackage{enumitem}
\usepackage[most]{tcolorbox}
\usepackage{xcolor}
\usepackage{fancyhdr}
\usepackage{pdflscape}
\usepackage[colorlinks=true,linkcolor=blue!50!black,urlcolor=blue!50!black]{hyperref}

\begin{document}

\begin{titlepage}
    \centering
    
    \vfill
    {\huge \textbf{PROCESAMIENTO DEL LENGUAJE NATURAL} \par}
    \vspace{2cm}
    {\Large Práctica 1: Preparación y procesamiento de textos \par}
    \vspace{1cm}
    {\Large Universidad Autónoma de Madrid \par}
    \vspace{1cm}
    {\large Iván Domínguez Hernández |  ivan.dominguezh@estudiante.uam.es \par}
    \vspace{.5cm}
    {\large Lucas Miranda López | lucas.mirandal@estudiante.uam.es \par}
    \vspace{1cm}
\end{titlepage}


\tableofcontents

\newpage

\section{Descarga y preparación de la información}
\label{sec-1}

\noindent \textbf{Resumen:} En esta primera sección se descarga de uno los ficheros del corpus (o base de datos textuales) 
de comentarios de \textit{Amazon} (uno sencillo y poco pesado, el de tarjetas de regalo) y se diseña la clase 
\texttt{ProcesadorDeTexto} para almacenar cada review en un fichero propio en el directorio que le corresponda sin almacenar 
reviews duplciadas. 

\begin{enumerate}
    \item La decisión de descargar el fichero \textit{Gift\_Cards.jsonl} se debe únicamente a su tamaño. 
    Teniendo en cuenta que todos ficheros del corpus de \textit{Amazon} siguen la misma estructura JSON, 
    lo óptimo para la realización de la práctica parecía ser escoger un fichero ligero sobre el que hacer
    pruebas en un archivo dde tipo \textit{.ipynb} para, posteriormente, codificar sencillamente el archivo 
    \textit{.py} definitvo y genérico capaz de procesar cualquier fichero del corpus. 

    \item Las pruebas realizadas sobre un archivo \textit{.ipynb} auxiliar, de prueba, y que no se entrega adjunto,
     permitieron construir la clase \texttt{ProcesadorDeTexto}, capaz de:

    \begin{enumerate}
        \item Recibir el nombre del fichero a preprocesar (sin especificar su extensión, solo el nombre del fichero). En 
        nuestro caso 'Gift\_Cards'.
        \item Recibir una lista con las categorías que el usuario quiera almacenar para cada reseña. Nosotros aconsejamos 
        pasarle la lista \texttt{['rating', 'title', 'text', 'parent\_asin', 'user\_id']}, que almacena la puntación en estrellas del 1 al 5, 
        el título (tal vez contenga información relevante), el texto per se de la reseña, una referencia al tipo de producto reseñado 
        (que posteriormente permitirá acceder a metainformación de ser preciso) y una referencia al usuario que hace la reseña. 
        \textbf{Nota:} la característica \texttt{'user\_id'} siempre es necesaria en la lista para asegurar que se pueden borrar los duplicados, 
        como se explica en el siguiente punto. Una vez usada, se elimina.\\
        
        Y ahora, iterando sobre cada elemento del \textit{.jsonl} a procesar:\\

        \item No almacenar reseñas que, siendo del mismo usuario y al mismo producto, tienen el mismo texto y título. Se puede pensar que 
        esto debería ser parte de un posterior preprocesado. Pero al implicar iterar sobre cada elemento de nuevo, decidimos hacerlo en este
        paso para ahorrar cómputo a posteriori. 
        \item Guardar las categorías enunciadas en la lista que se le pase de cada reseña en un fichero propio, en tres distintos directorios: 
        \textit{/positivas, /neutras} y \textit{/negativas}, a modo de etiqueta, basándonos únicamente a cual de los intervalos pertenece la
        categoría \textit{'rating'}: $rating\geq4.5$, $4.5 \geq rating \geq 1.5$ o $rating\leq1.5$. Es un criterio absurdo, 
        pero al ser configurable, de momento no nos preocuparemos por esto en esta práctica.  
    \end{enumerate}
 
\end{enumerate}

\noindent Por tanto, lo que podemos obtener tras esta sección es una división en directorios, basada en la puntuación de cada reseña, 
en los que se incluyen ficheros \textit{.json}, para cada reseña, y que contienen las categorías especificadas por el usuario. Esto
facilitará en gran medida la posterior ingesta (\ref{sec-2}) de datos por un gestor de corpus especializado en nuestro caso, 
\texttt{AmazonCorpusReader}. \textbf{Nota:} los directorios que se crean están bastante desbalanceados. Esto es, existen muchas más 
reseñas de clase positiva que de clases neutra o negativa. Esto es, en realidad, una distribución inherente, no solo al criterio 
seguido para el etiquetado, sino también a los propios datos. Y es que lo normal es que los usuarios hagan tantas reseñas negativas 
o neutras como positivas. 

\newpage

\section{Diseñar e implementar el gestor del corpus}
\label{sec-2}

\noindent \textbf{Resumen:} Para poder acceder al corpus en diferentes niveles (documento, párrafo, oración y 
palabra como indica el enunciado), así como a la categoría de cada reseña, se desarrolla la clase 
\texttt{AmazonCorpusReader}, heredera de \texttt{CategorizedCorpusReader} y \texttt{CorpusReader}. 
La estructura de la clase está fuertemente inspirada en \texttt{HTMLCorpusReader} propuesta en el libro \cite{bengfort2018}.


\begin{enumerate}
    \item \textbf{Herencia:} en el constructor (\texttt{\_\_init\_\_}) se inicializan ambas clases base por separado:
    \begin{enumerate}
        \item \texttt{CategorizedCorpusReader.\_\_init\_\_} recibe un \texttt{cat\_pattern}, una expresión regular 
        (\texttt{r'(positivas|neutras|negativas)/.*'}) que permite deducir automáticamente la categoría (\texttt{categories}) 
        de cada fichero a partir del directorio en el que se encuentra.
        \item \texttt{CorpusReader.\_\_init\_\_} recibe la ruta raíz del corpus y el patrón de los ficheros a considerar 
        (\texttt{fileids}, por defecto \texttt{r'.*\textbackslash.json'}).
    \end{enumerate}

    \item \textbf{Método \texttt{\_resolve()}}: dado que la mayoría de métodos de la clase pueden recibir una lista de 
    \texttt{fileids} o una lista de \texttt{categories}, este método auxiliar (inspirado en el patrón que sigue 
    \texttt{CategorizedPlaintextCorpusReader} en NLTK) unifica ambos casos: si se reciben categorías, las traduce a 
    la lista de ficheros correspondiente mediante \texttt{self.fileids(categories)}; si se reciben ambas a la vez, 
    lanza un \texttt{ValueError}, al no tener sentido combinarlas.

    \item \textbf{Método \texttt{docs()}}: recorre uno a uno los ficheros seleccionados (vía \texttt{\_resolve()}) y, 
    para cada línea del \texttt{.json} correspondiente, construye un diccionario con sus campos (\texttt{title}, 
    \texttt{text}, \texttt{rating}, \texttt{parent\_asin}\ldots) y le añade la clave \texttt{'category'} obtenida 
    mediante el método \texttt{\_resolve\_with\_category()}. El resto de métodos de acceso hacen uso de él.

    \item \textbf{Métodos de acceso a campos concretos:} a partir de \texttt{docs()} se construyen, como generadores, 
    \texttt{titles()}, \texttt{ratings()} y \texttt{texts()}, que devuelven respectivamente cada uno el título, la puntuación 
    y el texto de cada reseña seleccionada, y \texttt{raw()}, que devuelve el contenido en crudo de los ficheros unidos, 
    reutilizando la función \texttt{concat()} de NLTK.

    \item \textbf{Métodos de segmentación:} implementa la jerarquía habitual de un corpus reader de NLTK.
    \begin{enumerate}
        \item \texttt{paras()}: para cada reseña, separa su texto en párrafos (\texttt{BlanklineTokenizer}), cada párrafo en 
        oraciones (\texttt{sent\_tokenize}) y cada oración en palabras (\texttt{word\_tokenize}), devolviendo una estructura 
        anidada lista $\to$ lista $\to$ lista.
        \item \texttt{sents()}: itera sobre \texttt{paras()} y «desenvuelve» un nivel, devolviendo cada oración (como lista de 
        palabras) de forma independiente.
        \item \texttt{words()}: itera a su vez sobre \texttt{sents()} y devuelve cada palabra individual del corpus seleccionado.
    \end{enumerate}
    Todos estos métodos están implementados como generadores (\texttt{yield}), de forma que el corpus nunca se carga por completo 
    en memoria, sino que se recorre reseña a reseña.

    \item \textbf{Método \texttt{describe()}:} realiza una única pasada sobre el corpus (o, más bien, el subconjunto de 
    \texttt{fileids}/\texttt{categories} indicado) y devuelve un diccionario con las métricas habituales de caracterización 
    de corpus: número de ficheros (\texttt{files}) y de categorías (\texttt{topics}) consideradas; número total de párrafos, 
    oraciones y palabras (\texttt{paras}, \texttt{sents}, \texttt{words}); tamaño del vocabulario (\texttt{vocab}); diversidad 
    léxica (\texttt{lexdiv}, palabras entre vocabulario); párrafos por documento (\texttt{ppdoc}); oraciones por párrafo 
    (\texttt{sppar}); y tiempo de cómputo empleado (\texttt{secs}). Para ello utiliza dos objetos \texttt{nltk.FreqDist()} 
    (uno para los contadores globales y otro para la frecuencia de cada token), evitando recorrer el corpus más de una vez.
\end{enumerate}

\newpage

\section{Preprocesamiento del corpus y análisis básicos}
\label{sec-3}

\noindent El preprocesamiento se incorpora a nivel de corpus reader. Para hacerlo, fue necesario (dejando de lado la actualización 
de NLTK), si bien no modificar la funcionalidad ya implementada en la sección \ref{sec-2}, añadir bloques lógicos en algunos puntos
para la propia implementación de la nueva funcionalidad de preprocesado. En particular, los cambios fueron los siguientes:

\begin{enumerate}
    \item Creación del método \texttt{\_limpiar()}: limpia el texto de acuerdo con los requisitos del enunciado de la práctica, en 
    el siguiente orden: \begin{enumerate}
        \item Eliminar etiquetas HTML, de haberlas. Todas pasan a espacio en blanco, salvo \texttt{<br />} que pasa a ser salto de 
        línea. 
        \item Convierte entidades \texttt{\&amp} $\to$ \texttt{\&}. 
        \item Elimina emojis. 
        \item Reduce puntuación repetida !!! $\to$ !. 
        \item Reduce letras repetidas \textit{gooooood} $\to$ \textit{good}. 
        \item Normaliza espacios conservando saltos de línea. 
    \end{enumerate}
    Con tan solo añadir la línea \texttt{obj['text'] = self.\_limpiar(obj['text'])}, ahora todo el resto de métodos 
    (\texttt{texts()}, \texttt{paras()}, \texttt{sents()}, \texttt{words()} y \texttt{describe()}) reciben el texto ya limpio. 
    \item El método \texttt{paras()} se rehace para evitar blancos: el método ahora comprueba en cada nivel 
    (\texttt{para.strip()}, \texttt{sent.strip()}, \texttt{w.strip()}) y si es blanco se descarta (manteniendo el \texttt{yield}).
    \item Método \texttt{\_get\_stopwords}: los métodos \texttt{paras()}, \texttt{sents()} y \texttt{words()} reciben ahora
    un parámetro \texttt{sin\_stopwords = False}. Con \texttt{False}, se eliminan las que NLTK asume como stopwords (en inglés) 
    a excepción de 'no', 'not' y 'nor' para evitar cambios de significado, con \texttt{True}, no. 
    \item Retoques en el método \texttt{describe()}: ahora tiene en cuenta el número total de caracteres (se nos olvidó hacerlo)
    en la sección \ref{sec-2}, además de aceptar el parámetro \texttt{sin\_stepwords}. 
\end{enumerate}

\section{Estadísticas del corpus}
\label{sec-4}

\noindent \textbf{Resumen:} Para analizar el corpus se crea el script \textit{estadisticas.py}, que no implementa
ninguna métrica desde cero, sino que aprovecha los métodos ya implementados en \texttt{AmazonCorpusReader}
(secciones \ref{sec-2} y \ref{sec-3}). La base del análisis es el método \texttt{describe()}, al que se le suman
algunos análisis que este no ofrece.

\begin{enumerate}
    \item \textbf{Métodos de la clase utilizados:}
    \begin{enumerate}
        \item \texttt{categories()} y \texttt{fileids()}: número de reseñas por categoría, es decir, el balance de clases.
        \item \texttt{describe()}: métricas generales de cada categoría. Se llama dos veces, con y sin \textit{stopwords},
        para ver cuánto texto suponen estas.
        \item \texttt{ratings()}: distribución de estrellas de cada categoría. Sirve, además, para comprobar que el
        etiquetado de la sección \ref{sec-1} se hizo bien.
        \item \texttt{titles()} y \texttt{paras()}: longitud de títulos y reseñas (palabras y oraciones por reseña, mediana,
        máximo y reseñas que quedan vacías tras la limpieza).
        \item \texttt{words(sin\_stopwords=True)}: palabras más frecuentes de cada categoría, sin signos de puntuación.
        A partir de estas frecuencias se obtienen también las palabras \textit{características} de cada categoría:
        las que tienen mayor frecuencia relativa en ella que en el resto.
    \end{enumerate}
    \item \textbf{Ejecución:} \texttt{python3 estadisticas.py} analiza el corpus completo (unos 2 minutos). Con
    \texttt{-{}-muestra N} se analiza un subconjunto aleatorio (con semilla fija, para que sea reproducible) de, como
    mucho, $N$ reseñas por categoría. Los resultados se guardan en \textit{estadisticas.json} y en \textit{estadisticas.tex},
    que se incluye directamente en esta memoria (tablas \ref{tab:estadisticas} y \ref{tab:top-palabras}).
\end{enumerate}

\begin{table}[h]
\centering
\begin{tabular}{lrrr}
\toprule
Métrica & negativas & neutras & positivas \\
\midrule
Reseñas en el corpus & 12.100 & 11.568 & 124.439 \\
\% del corpus & 8,17 & 7,81 & 84,02 \\
Reseñas analizadas & 12.100 & 11.568 & 124.439 \\
Caracteres & 2.418.048 & 1.584.797 & 8.327.376 \\
Párrafos & 12.917 & 12.113 & 125.658 \\
Oraciones & 36.934 & 25.572 & 196.146 \\
Palabras & 526.921 & 346.770 & 1.821.590 \\
Palabras (sin stopwords) & 299.204 & 195.276 & 1.095.203 \\
Vocabulario & 12.749 & 10.648 & 23.702 \\
\texttt{lexdiv} & 41,33 & 32,57 & 76,85 \\
Párrafos por reseña & 1,07 & 1,05 & 1,01 \\
Oraciones por párrafo & 2,86 & 2,11 & 1,56 \\
Palabras por reseña & 43,55 & 29,98 & 14,64 \\
Rating medio & 1,00 & 3,40 & 5,00 \\
\bottomrule
\end{tabular}
\caption{Estadísticas del corpus por categoría (\texttt{AmazonCorpusReader}).}
\label{tab:estadisticas}
\end{table}

\begin{table}[h]
\centering
\begin{tabularx}{\textwidth}{lX}
\toprule
Categoría & Palabras más frecuentes (sin stopwords) \\
\midrule
negativas & card, gift, not, amazon, cards, n't, no, use, get, money \\
neutras & gift, card, not, amazon, cards, n't, use, one, no, would \\
positivas & gift, card, great, amazon, easy, cards, use, love, perfect, good \\
\bottomrule
\end{tabularx}
\caption{Palabras más frecuentes por categoría.}
\label{tab:top-palabras}
\end{table}

\noindent De las tablas anteriores (corpus completo) se extrae lo siguiente:

\begin{enumerate}
    \item \textbf{Desbalanceo:} se confirma lo que ya se intuía en la sección \ref{sec-1}: el 84\,\% de las reseñas son
    positivas, frente a un 8,2\,\% de negativas y un 7,8\,\% de neutras.
    \item \textbf{Longitud:} cuanto peor es la reseña, más larga. Una negativa tiene de media unas 44 palabras y 3 oraciones,
    una neutra unas 30 y una positiva apenas 15 (el 66,5\,\% de las positivas tienen una única oración). Tiene sentido:
    quien está contento escribe «Great gift» y se va, pero quien ha tenido un problema lo cuenta con detalle.
    \item \textbf{\textit{Stopwords}:} suponen alrededor del 40-44\,\% de las palabras del corpus, por lo que eliminarlas
    reduce casi a la mitad el texto a procesar.
    \item \textbf{Vocabulario:} las palabras más frecuentes (\textit{gift}, \textit{card}, \textit{amazon}) son comunes a
    las tres categorías y no sirven para distinguirlas, pues son propias del producto. Las palabras características,
    en cambio, sí son muy reveladoras: \textit{dispute}, \textit{scam} o \textit{fraudulent} en las negativas;
    \textit{tacky}, \textit{dented} o \textit{slow} en las neutras; y \textit{cutest}, \textit{no-brainer} o
    \textit{perfect} en las positivas. Además, \textit{not} y \textit{n't} aparecen entre las más frecuentes de negativas
    y neutras, lo que justifica haberlas conservado al filtrar las \textit{stopwords} (sección \ref{sec-3}).
    \item \textbf{Otros idiomas y reseñas vacías:} entre las palabras características de las positivas aparecen
    \textit{excelente} y \textit{perfecto}, es decir, hay reseñas en español (ver comentarios). Por otro lado, 397 reseñas
    (351 de ellas positivas) quedan vacías tras la limpieza, al contener únicamente emojis o etiquetas HTML.
\end{enumerate}

\noindent \textbf{Nota:} \texttt{lexdiv} (y también \texttt{vocab}) depende del tamaño del texto: cuantas más palabras,
más se repite cada una. Por eso las positivas tienen un \texttt{lexdiv} mucho mayor (76,85 frente a 41,33 de las
negativas), pero no porque su vocabulario sea más pobre, sino porque tienen casi 3,5 veces más palabras que las
negativas. Para comparar estas métricas entre categorías de forma justa conviene usar
\texttt{-{}-muestra}, de forma que las tres tengan el mismo número de reseñas.

\section{Comentarios y conclusiones}

\subsection{Comentarios}

\begin{enumerate}
    \item Tal vez el hecho de eliminar repetidos no fuera una tarea de la sección \ref{sec-1}, pues consiste más bien en 
    preprocesamiento. Sin embargo, se nos ocurrió hacerlo el primer día de práctica, sin mucha visión de lo que sería el 
    resultado final de la misma, por lo que decidimos dejarlo en esa sección. Además, se gana en términos de eficiencia, pues 
    ya no es necesario volver a iterar sobre cada uno de los comentarios del corpus. 
    \item Como propuesta exploratoria para otras prácticas, es probable que convenga integrar la clase \texttt{ProcesadorDeTexto}
    con herramientas de procesamiento por lotes y distribuido, pues es una tarea muy pesada y paralelizable. Un comentario 
    también relacionado con esta clase es que el nombre escogido no es el mejor, pues no es un procesador ni tampoco procesa texto
    como tal. Su tarea es más bien darle una estructura de corpus a los datos ofrecidos por \textit{Amazon}. 
    \item No teníamos muy claro si era lo suyo eliminar todos los comentarios escritos en un idioma distinto del inglés. 
    Nuestra teoría es que si no hay mucho desbalanceo de comentarios entre texto escrito en inglés y texto escrito en otros 
    idiomas, sabiendo que la cantidad de comentarios no es un problema, un modelo terminaría por aprender a clasificar sin
    fijarse en el lenguaje en el que está escrito. 
    \item Tuvimos problemas con la versión de NLTK porque aparentemente las versiones que manejábamos eran distintas y 
    la funcionalidad cambiaba entre una y otra. Por favor, de encontrar algún problema en cuanto a la ejecución del 
    código, nótese que puede ser debido a esto. Antes de suspendernos, consúltanos :`(
    \item La implementación de todos los códigos (salvo el de pruebas y casos de uso, \textit{prueba.py}) de esta práctica 
    es nuestra, si bien nos hemos ayudado en herramientas de inteligencia artificial generativa (Claude Sonnet 5.5)
    (sección \ref{sec-3}), así como en los recursos bibliográficos propuestos en la práctica (sección \ref{sec-2}, \cite{bengfort2018}). 
\end{enumerate}

\subsection{Conclusiones}

\begin{enumerate}
    \item El corpus está desbalanceado hacia reseñas positivas. Esto es algo completamente irrelevante para esta 
    primera práctica pero asumimos que no lo será para prácticas donde entrenemos modelos. Por ello, como hipótesis, 
    creemos que la mejor forma de tratar esto sería el \textit{undersampling} para ahorrar cómputo. 
    \item Las \textit{stopwords} suponen una gran parte del texto, pero no apotan información especialmente relevante 
    acerca de si la reseña es positiva, neutra o negativa (por eso se han eliminado todas a excepción de 'no', 'nor' y 
    'not'). 

\end{enumerate}

\newpage

\begin{thebibliography}{9}
\bibitem{bengfort2018}
B. Bengfort, R. Bilbro y T. Ojeda,
\textit{Applied Text Analysis with Python: Enabling Language-Aware Data Products with Machine Learning},
1.ª ed., O'Reilly Media, Sebastopol, CA, 2018. ISBN 978-1-4919-6304-3.

\end{thebibliography}

\end{document}